\section{Related Work}

The closest systems neighbor is CodecSight \cite{codecsight}, which also tries
to avoid unnecessary visual processing by using compressed-video structure
before the rest of the model consumes the content. The difference is both in
signal and in target. CodecSight uses decoder-native motion vectors and
residuals and is evaluated on streaming anomaly detection; our current local
implementation uses a pixel-domain proxy because that is what our MLX stack can
exercise today, and our main benchmark evaluations target temporal-reasoning QA
on TOMATO \cite{tomato} and MVBench \cite{mvbench}.

Trained codec-native approaches answer a different question. CoPE-VideoLM
\cite{copevideolm}, Deja Vu \cite{dejavu}, and older compressed-video pipelines
such as CoViAR \cite{coviar} learn how motion, residual, or reused
representations should enter the model. That is valuable prior art, but it does
not settle the deployment question we care about here: how much redundancy a
frozen stack can exploit before retraining or representation learning enters
the picture.

Another family reduces visual cost after tokens already exist. FastV
\cite{fastv}, FastVID \cite{fastvid}, FlashVID \cite{flashvid}, VisionZip
\cite{visionzip}, SparseVLM \cite{sparsevlm}, VScan \cite{vscan}, and STTM
\cite{sttm} prune or merge visual tokens using attention-derived, density, or
feature-similarity signals. Those methods are adjacent but not identical to our
setting. Their signal usually lives after visual tokens exist. Our routing lane
asks an earlier question: can we decide where fresh computation should land in
time before the full dense backend pays for every frame again?

Long-horizon reuse and systems co-design are also relevant. QuickVideo
\cite{quickvideo} explicitly co-designs decode and prefill. ReKV \cite{rekv},
VLCache \cite{vlcache}, and StreamingVLM \cite{streamingvlm} study how cached
visual or KV state can be reused across requests or over long contexts. That
line of work is especially relevant to our persistent-KV contribution: the
paper's largest local speedups do not come from a better frame selector, but
from refusing to repay the same prefix when the next question is about the same
video.

SimpleStream \cite{simplestream} is also an important streaming baseline
because it argues that a strong recent-frame window can match or outperform
heavier streaming-memory systems on modern benchmarks. That critique matters
directly for our streaming lane: any cache or memory mechanism has to beat a
strong recency baseline under matched budget, not just beat a strawman.

More conceptually, our framing is adjacent to approximate computing
\cite{rinardapprox}: skipping work under an explicit quality budget. The
difference is that the skipped work here is temporal visual recomputation, and
the guardrail is benchmark fidelity rather than an application-specific numeric
tolerance.

Our position is therefore narrower and more explicit than a generic
``efficient-video-VLM'' claim. We study training-free anti-recomputation for
frozen video VLMs, separating first-pass reuse, after-ingest reuse, and
mechanism-validation routing. The paper is not a claim that pixel differencing
is the final signal. It is a claim that careful temporal reuse is already worth
studying before the full codec-native path lands.

\begin{sectiontodo}{Freeze the final comparison table against the primary papers}
The mechanism families are now in the right buckets, but a submission version
still needs a tighter final comparison table with exact titles, venues, and one
checked headline metric per cited method.

\textcolor{todoBorder}{Source pointers: \path{docs/related-work-table.md},
\path{docs/literature-map-2026-04-16.md},
\path{paper/framing.md}.}
\end{sectiontodo}
